% GENERATED FILE. Edit canonical lessons, then run npm run generate:books.
% canonical-source-hash: fnv1a64:83ecf0cbbbcdadb5
% canonical-lessons: GU-C189-lapetvu, GU-C189-mangavvu, GU-C189-khutvu, GU-C189-bacvu, GU-C189-svikarvu

\chapter{Wrapping up an order}
\label{ch:gu-wrapping-up-an-order}
\hlchaptermodality{189}

\begin{chapteropening}
\textbf{By the end of this chapter:} \emph{I can say what is wrapped, ordered, used up or left over.}

Complete the last lesson of chapter 189: I can say what is wrapped, ordered, used up or left over.
\end{chapteropening}

\section[\texorpdfstring{\gu{લપેટવું} (lapeṭvũ)}{લપેટવું (lapeṭvũ)}]{\gu{લપેટવું} (lapeṭvũ) — to wrap}

\label{lesson:GU-C189-lapetvu}



\begin{quote}
Before the new one: say the Gujarati for goods, then the Gujarati for thing.
\end{quote}

\subsection*{You'll want to know: \gu{લપેટવું}}
\textbf{\gu{લપેટવું}} — \emph{lapeṭvũ} — \textquotedblleft{}to wrap\textquotedblright{}.

\begin{quote}
\textbf{\gu{દુકાનદાર} \gu{ભેટ} \gu{કાગળમાં} \gu{લપેટે} \gu{છે}.} — \emph{dukāndār bheṭ kāgaḷmā̃ lapeṭe che} — \textquotedblleft{}The shopkeeper wraps the gift in paper.\textquotedblright{}
\end{quote}

\begin{cousinweb}
\textbf{\gu{લપેટવું}} (\emph{lapeṭvũ}) is a tadbhava. It comes from \emph{*lappēṭṭ}, \textquotedblleft{}wrap\textquotedblright{}, a reconstructed form, and CDIAL glosses the Gujarati verb \textquotedblleft{}to wrap, fold, wind\textquotedblright{} (CDIAL 10942).
\end{cousinweb}

\begin{grammarlens}[title={What you've built}]
One more word for everyday life at A2.
\end{grammarlens}

\subsection*{Your turn}
\begin{itemize}
\raggedright
  \item \emph{Say it:} \emph{lapeṭvũ}
  \item \emph{Say it:} \emph{lapeṭvũ}, once more
  \item \emph{Say it:} \emph{dukāndār bheṭ kāgaḷmā̃ lapeṭe che} — The shopkeeper wraps the gift in paper.
  \item \emph{Recall:} say the Gujarati for goods, then the Gujarati for thing, then say \emph{lapeṭvũ} again
\end{itemize}

\begin{tcolorbox}[breakable,colback=teal!4,colframe=teal!35!black,title={Before you move on}]
What does \gu{લપેટવું} mean? (\textquotedblleft{}To wrap\textquotedblright{}.) Say it once more.
\end{tcolorbox}
\section[\texorpdfstring{\gu{મંગાવવું} (mangāvvũ)}{મંગાવવું (mangāvvũ)}]{\gu{મંગાવવું} (mangāvvũ) — to order, to send for}

\label{lesson:GU-C189-mangavvu}



\begin{quote}
Before the new one: say the Gujarati for thing, then the Gujarati for to wrap.
\end{quote}

\subsection*{You'll want to know: \gu{મંગાવવું}}
\textbf{\gu{મંગાવવું}} — \emph{mangāvvũ} — \textquotedblleft{}to order, to send for\textquotedblright{}.

\begin{quote}
\textbf{\gu{અમે} \gu{હોટેલમાંથી} \gu{ચા} \gu{મંગાવીએ} \gu{છીએ}.} — \emph{ame hoṭelmā̃thī cā mangāvīe chīe} — \textquotedblleft{}We order tea from the hotel.\textquotedblright{}
\end{quote}

\textbf{\gu{માગવું}} asks for a thing; \textbf{\gu{મંગાવવું}} has it brought.

\begin{cousinweb}
\textbf{\gu{મંગાવવું}} (\emph{mangāvvũ}) is built on \emph{māgvũ}, \textquotedblleft{}to ask for\textquotedblright{}, with the causative ending \emph{-āvvũ}: \textquotedblleft{}to have something brought\textquotedblright{}. \emph{māgvũ} is a tadbhava from Sanskrit \emph{m\'{\={a}}rgati}, \textquotedblleft{}seeks, requests\textquotedblright{} (CDIAL 10074).
\end{cousinweb}

\begin{grammarlens}[title={What you've built}]
One more word for everyday life at A2.
\end{grammarlens}

\subsection*{Your turn}
\begin{itemize}
\raggedright
  \item \emph{Say it:} \emph{mangāvvũ}
  \item \emph{Say it:} \emph{mangāvvũ}, once more
  \item \emph{Say it:} \emph{ame hoṭelmā̃thī cā mangāvīe chīe} — We order tea from the hotel.
  \item \emph{Recall:} say the Gujarati for thing, then the Gujarati for to wrap, then say \emph{mangāvvũ} again
\end{itemize}

\begin{tcolorbox}[breakable,colback=teal!4,colframe=teal!35!black,title={Before you move on}]
What does \gu{મંગાવવું} mean? (\textquotedblleft{}To order, to send for\textquotedblright{}.) Say it once more.
\end{tcolorbox}
\section[\texorpdfstring{\gu{ખૂટવું} (khūṭvũ)}{ખૂટવું (khūṭvũ)}]{\gu{ખૂટવું} (khūṭvũ) — to run out}

\label{lesson:GU-C189-khutvu}



\begin{quote}
Before the new one: say the Gujarati for to wrap, then the Gujarati for to order.
\end{quote}

\subsection*{You'll want to know: \gu{ખૂટવું}}
\textbf{\gu{ખૂટવું}} — \emph{khūṭvũ} — \textquotedblleft{}to run out\textquotedblright{}.

\begin{quote}
\textbf{\gu{દુકાનમાં} \gu{દૂધ} \gu{ખૂટ્યું}.} — \emph{dukānmā̃ dūdh khūṭyũ} — \textquotedblleft{}The shop has run out of milk.\textquotedblright{}
\end{quote}

\begin{cousinweb}
\textbf{\gu{ખૂટવું}} (\emph{khūṭvũ}) is a tadbhava. CDIAL lists it, glossed \textquotedblleft{}to be exhausted\textquotedblright{}, under \emph{*khuṭṭa}, \textquotedblleft{}defective\textquotedblright{}, a reconstructed form (CDIAL 3895).
\end{cousinweb}

\begin{grammarlens}[title={What you've built}]
One more word for everyday life at A2.
\end{grammarlens}

\subsection*{Your turn}
\begin{itemize}
\raggedright
  \item \emph{Say it:} \emph{khūṭvũ}
  \item \emph{Say it:} \emph{khūṭvũ}, once more
  \item \emph{Say it:} \emph{dukānmā̃ dūdh khūṭyũ} — The shop has run out of milk.
  \item \emph{Recall:} say the Gujarati for to wrap, then the Gujarati for to order, then say \emph{khūṭvũ} again
\end{itemize}

\begin{tcolorbox}[breakable,colback=teal!4,colframe=teal!35!black,title={Before you move on}]
What does \gu{ખૂટવું} mean? (\textquotedblleft{}To run out\textquotedblright{}.) Say it once more.
\end{tcolorbox}
\section[\texorpdfstring{\gu{બચવું} (bacvũ)}{બચવું (bacvũ)}]{\gu{બચવું} (bacvũ) — to be left over, to be spared}

\label{lesson:GU-C189-bacvu}



\begin{quote}
Before the new one: say the Gujarati for to order, then the Gujarati for to run out.
\end{quote}

\subsection*{You'll want to know: \gu{બચવું}}
\textbf{\gu{બચવું}} — \emph{bacvũ} — \textquotedblleft{}to be left over, to be spared\textquotedblright{}.

\begin{quote}
\textbf{\gu{મારી} \gu{પાસે} \gu{દસ} \gu{રૂપિયા} \gu{બચે} \gu{છે}.} — \emph{mārī pāse das rūpiyā bace che} — \textquotedblleft{}I have ten rupees left over.\textquotedblright{}
\end{quote}

\textbf{\gu{બચાવવું}} saves a thing; \textbf{\gu{બચવું}} is what is left or spared.

\begin{cousinweb}
\textbf{\gu{બચવું}} (\emph{bacvũ}) goes back to Sanskrit \emph{váñcati}, \textquotedblleft{}moves; escapes from\textquotedblright{}, but CDIAL marks the Gujarati verb, glossed \textquotedblleft{}to escape\textquotedblright{}, as borrowed from Hindi \emph{bacnā}, \textquotedblleft{}to escape\textquotedblright{}; the taught \emph{bacāvvũ}, \textquotedblleft{}to save\textquotedblright{}, is its causative (CDIAL 11208).
\end{cousinweb}

\begin{grammarlens}[title={What you've built}]
One more word for everyday life at A2.
\end{grammarlens}

\subsection*{Your turn}
\begin{itemize}
\raggedright
  \item \emph{Say it:} \emph{bacvũ}
  \item \emph{Say it:} \emph{bacvũ}, once more
  \item \emph{Say it:} \emph{mārī pāse das rūpiyā bace che} — I have ten rupees left over.
  \item \emph{Recall:} say the Gujarati for to order, then the Gujarati for to run out, then say \emph{bacvũ} again
\end{itemize}

\begin{tcolorbox}[breakable,colback=teal!4,colframe=teal!35!black,title={Before you move on}]
What does \gu{બચવું} mean? (\textquotedblleft{}To be left over, to be spared\textquotedblright{}.) Say it once more.
\end{tcolorbox}
\section[\texorpdfstring{\gu{સ્વીકારવું} (svīkārvũ)}{સ્વીકારવું (svīkārvũ)}]{\gu{સ્વીકારવું} (svīkārvũ) — to accept}

\label{lesson:GU-C189-svikarvu}



\begin{quote}
Before the new one: say the Gujarati for to run out, then the Gujarati for to be left over.
\end{quote}

\subsection*{You'll want to know: \gu{સ્વીકારવું}}
\textbf{\gu{સ્વીકારવું}} — \emph{svīkārvũ} — \textquotedblleft{}to accept\textquotedblright{}.

\begin{quote}
\textbf{\gu{આ} \gu{દુકાન} \gu{સિક્કા} \gu{સ્વીકારે} \gu{છે}.} — \emph{ā dukān sikkā svīkāre che} — \textquotedblleft{}This shop accepts coins.\textquotedblright{}
\end{quote}

\begin{cousinweb}
\textbf{\gu{સ્વીકારવું}} (\emph{svīkārvũ}) is built on Sanskrit \emph{svīkāra}, \textquotedblleft{}acceptance, making one's own\textquotedblright{}, from \emph{sva}, \textquotedblleft{}own\textquotedblright{}, and \emph{kāra}, \textquotedblleft{}making\textquotedblright{}: a learned word with the Gujarati verb ending \emph{-vũ}. CDIAL has no entry for it.
\end{cousinweb}

\begin{grammarlens}[title={What you've built}]
One more word for everyday life at A2.
\end{grammarlens}

\subsection*{Your turn}
\begin{itemize}
\raggedright
  \item \emph{Say it:} \emph{svīkārvũ}
  \item \emph{Say it:} \emph{svīkārvũ}, once more
  \item \emph{Say it:} \emph{ā dukān sikkā svīkāre che} — This shop accepts coins.
  \item \emph{Recall:} say the Gujarati for to run out, then the Gujarati for to be left over, then say \emph{svīkārvũ} again
\end{itemize}

\begin{tcolorbox}[breakable,colback=teal!4,colframe=teal!35!black,title={Before you move on}]
What does \gu{સ્વીકારવું} mean? (\textquotedblleft{}To accept\textquotedblright{}.) Say it once more.
\end{tcolorbox}
